\documentclass[preview]{standalone}

\usepackage{amsmath}
\usepackage{amssymb}
\usepackage{stellar}
\usepackage{definitions}
\usepackage{bettelini}

\begin{document}

\id{probability-limit-theorems}
\genpage

\section{Sample means}

\begin{snippetdefinition}{iid-random-variables-definition}{Independent and identically distributed random variables}
    Let \((\Omega,\mathcal{F},\mathbb{P})\) be a \probabilityspace, and
    let \((X_n)_{n\geq1}\) be a sequence of real-valued
    \randomvariable[random variables] on \(\Omega\). The sequence
    \((X_n)_{n\geq1}\) is \emph{independent and identically distributed}
    if the \randomvariable[random variables] \(X_n\) are mutually
    independent and
    \[
        X_n\eqdist X_1
    \]
    for every \(n\geq1\).
\end{snippetdefinition}

\begin{snippetdefinition}{sample-mean-definition}{Sample mean}
    Let \(X_1,\cdots,X_n\) be real-valued \randomvariable[random variables]
    on a \probabilityspace. The \emph{sample mean} of
    \(X_1,\cdots,X_n\) is the \randomvariable
    \[
        \overline{X}_n
        =
        \frac{1}{n}\sum_{i=1}^{n}X_i.
    \]
\end{snippetdefinition}

\begin{snippetexample}{bernoulli-sample-mean-motivation}{Bernoulli sample means}
    Let \((X_n)_{n\geq1}\) be \iidrandomvariables with
    \[
        X_n\sim\operatorname{Bernoulli}(p)
    \]
    for some \(p\in[0,1]\). The \samplemean
    \[
        \overline{X}_n=\frac{1}{n}\sum_{i=1}^{n}X_i
    \]
    is the random relative frequency of successes among the first
    \(n\) trials. Limit theorems describe in which sense
    \(\overline{X}_n\) approaches the parameter \(p\) as \(n\to\infty\).
\end{snippetexample}

\section{Markov and Chebyshev inequalities}

\begin{snippettheorem}{markov-inequality}{Markov inequality}
    Let \(X\) be a non-negative \randomvariable on a \probabilityspace.
    If \(\expectation[X]\) exists and is finite, then, for every
    \(a>0\),
    \[
        \mathbb{P}(X\geq a)\leq\frac{\expectation[X]}{a}.
    \]
\end{snippettheorem}

\begin{snippetproof}{markov-inequality-proof}{markov-inequality}{Markov inequality}
    Since \(X\geq0\),
    \[
        X\geq a\indicator_{\{X\geq a\}}.
    \]
    Taking expectations gives
    \[
        \expectation[X]
        \geq
        a\,\expectation[\indicator_{\{X\geq a\}}]
        =
        a\,\mathbb{P}(X\geq a).
    \]
    Dividing by \(a>0\) yields the inequality.
\end{snippetproof}

\begin{snippettheorem}{chebyshev-inequality}{Chebyshev inequality}
    Let \(X\) be a real-valued \randomvariable on a \probabilityspace.
    If \(\expectation[X]\) and \(\variance(X)\) exist and
    \(\variance(X)<+\infty\), then, for every \(\varepsilon>0\),
    \[
        \mathbb{P}\left(|X-\expectation[X]|\geq\varepsilon\right)
        \leq
        \frac{\variance(X)}{\varepsilon^2}.
    \]
\end{snippettheorem}

\begin{snippetproof}{chebyshev-inequality-proof}{chebyshev-inequality}{Chebyshev inequality}
    Apply the \markovinequality to the non-negative \randomvariable
    \[
        Y=(X-\expectation[X])^2
    \]
    with \(a=\varepsilon^2\). Then
    \[
        \mathbb{P}\left(|X-\expectation[X]|\geq\varepsilon\right)
        =
        \mathbb{P}\left((X-\expectation[X])^2\geq\varepsilon^2\right)
        \leq
        \frac{\expectation[(X-\expectation[X])^2]}{\varepsilon^2}
        =
        \frac{\variance(X)}{\varepsilon^2}.
    \]
\end{snippetproof}

\section{Laws of large numbers}

\begin{snippettheorem}{weak-law-large-numbers}{Weak law of large numbers}
    Let \((X_n)_{n\geq1}\) be a sequence of independent real-valued
    \randomvariable[random variables] on a \probabilityspace. Assume that,
    for every \(n\geq1\),
    \[
        \expectation[X_n]=\mu
        \quad\text{and}\quad
        \variance(X_n)=\sigma^2<+\infty.
    \]
    Then
    \[
        \overline{X}_n\xrightarrow{\mathbb{P}}\mu,
    \]
    where
    \[
        \overline{X}_n=\frac{1}{n}\sum_{i=1}^{n}X_i.
    \]
\end{snippettheorem}

\begin{snippetproof}{weak-law-large-numbers-proof}{weak-law-large-numbers}{Weak law of large numbers}
    For every \(n\geq1\),
    \[
        \expectation[\overline{X}_n]
        =
        \frac{1}{n}\sum_{i=1}^{n}\expectation[X_i]
        =
        \mu.
    \]
    Since the \randomvariable[random variables] \(X_i\) are independent,
    \[
        \variance(\overline{X}_n)
        =
        \frac{1}{n^2}\sum_{i=1}^{n}\variance(X_i)
        =
        \frac{\sigma^2}{n}.
    \]
    By the \chebyshevinequality, for every \(\varepsilon>0\),
    \[
        \mathbb{P}\left(|\overline{X}_n-\mu|\geq\varepsilon\right)
        \leq
        \frac{\sigma^2}{n\varepsilon^2}.
    \]
    The right-hand side converges to \(0\), so
    \[
        \overline{X}_n\xrightarrow{\mathbb{P}}\mu.
    \]
\end{snippetproof}

\begin{snippetexample}{bernoulli-sample-mean-chebyshev-bound}{Chebyshev bound for Bernoulli sample means}
    Let \((X_n)_{n\geq1}\) be \iidrandomvariables with
    \(X_n\sim\operatorname{Bernoulli}(p)\). Then
    \[
        \expectation[X_n]=p
        \quad\text{and}\quad
        \variance(X_n)=p(1-p).
    \]
    For every \(\varepsilon>0\), the \chebyshevinequality gives
    \[
        \mathbb{P}\left(|\overline{X}_n-p|\geq\varepsilon\right)
        \leq
        \frac{p(1-p)}{n\varepsilon^2}
        \leq
        \frac{1}{4n\varepsilon^2}.
    \]
    Hence the random relative frequency of successes converges in
    probability to \(p\).
\end{snippetexample}

\begin{snippettheorem}{strong-law-large-numbers}{Strong law of large numbers}
    Let \((X_n)_{n\geq1}\) be a sequence of \iidrandomvariables on a
    \probabilityspace. Assume that
    \[
        \expectation[X_1]=\mu
        \quad\text{and}\quad
        \variance(X_1)<+\infty.
    \]
    Then
    \[
        \overline{X}_n\xrightarrow{\mathrm{a.s.}}\mu,
    \]
    where
    \[
        \overline{X}_n=\frac{1}{n}\sum_{i=1}^{n}X_i.
    \]
\end{snippettheorem}

\begin{snippetproof}{strong-law-large-numbers-proof}{strong-law-large-numbers}{Strong law of large numbers}
    We first prove the result under the additional assumption
    \(X_1\geq0\). Let
    \[
        S_n=\sum_{i=1}^{n}X_i.
    \]
    By the \chebyshevinequality, for every \(\varepsilon>0\),
    \[
        \mathbb{P}\left(\left|\frac{S_{k^2}}{k^2}-\mu\right|
        \geq\varepsilon\right)
        \leq
        \frac{\variance(X_1)}{k^2\varepsilon^2}.
    \]
    Therefore
    \[
        \sum_{k=1}^{+\infty}
        \mathbb{P}\left(\left|\frac{S_{k^2}}{k^2}-\mu\right|
        \geq\varepsilon\right)<+\infty.
    \]
    By the \borelcantellilemma,
    \[
        \frac{S_{k^2}}{k^2}\xrightarrow{\mathrm{a.s.}}\mu.
    \]

    Since \(X_i\geq0\), the sequence \((S_n)\) is increasing. If
    \(k^2\leq n<(k+1)^2\), then
    \[
        \frac{k^2}{(k+1)^2}\frac{S_{k^2}}{k^2}
        \leq
        \frac{S_n}{n}
        \leq
        \frac{(k+1)^2}{k^2}\frac{S_{(k+1)^2}}{(k+1)^2}.
    \]
    Both outer terms converge almost surely to \(\mu\), so
    \[
        \frac{S_n}{n}\xrightarrow{\mathrm{a.s.}}\mu
    \]
    in the non-negative case.

    For a general real-valued \(X_1\), write
    \[
        X_i=X_i^+-X_i^-,
        \quad
        X_i^+=\max\{X_i,0\},
        \quad
        X_i^-=\max\{-X_i,0\}.
    \]
    The sequences \((X_i^+)_{i\geq1}\) and \((X_i^-)_{i\geq1}\) are
    \iidrandomvariables, non-negative, and have finite variance. Applying
    the non-negative case to both sequences gives
    \[
        \frac{1}{n}\sum_{i=1}^{n}X_i^+
        \xrightarrow{\mathrm{a.s.}}
        \expectation[X_1^+],
        \quad
        \frac{1}{n}\sum_{i=1}^{n}X_i^-
        \xrightarrow{\mathrm{a.s.}}
        \expectation[X_1^-].
    \]
    Subtracting the two limits gives
    \[
        \overline{X}_n
        \xrightarrow{\mathrm{a.s.}}
        \expectation[X_1^+]-\expectation[X_1^-]
        =
        \expectation[X_1]
        =
        \mu.
    \]
\end{snippetproof}

\begin{snippetcorollary}{strong-law-implies-weak-law}{The strong law implies the weak law}
    Let \((X_n)_{n\geq1}\) be a sequence of \iidrandomvariables on a
    \probabilityspace. Assume that
    \[
        \expectation[X_1]=\mu
        \quad\text{and}\quad
        \variance(X_1)<+\infty.
    \]
    Then the \stronglawlargenumbers implies the \weaklawlargenumbers:
    \[
        \overline{X}_n\xrightarrow{\mathrm{a.s.}}\mu
        \quad\Longrightarrow\quad
        \overline{X}_n\xrightarrow{\mathbb{P}}\mu.
    \]
\end{snippetcorollary}

\begin{snippetproof}{strong-law-implies-weak-law-proof}{strong-law-implies-weak-law}{The strong law implies the weak law}
    This is the implication from \almostsureconvergence to
    \convergenceinprobability applied to the sequence
    \((\overline{X}_n)_{n\geq1}\) and the constant \randomvariable equal
    to \(\mu\).
\end{snippetproof}

\section{Central limit theorem}

\begin{snippettheorem}{central-limit-theorem}{Central limit theorem}
    Let \((X_n)_{n\geq1}\) be a sequence of \iidrandomvariables on a
    \probabilityspace. Assume that
    \[
        \expectation[X_1]=\mu
        \quad\text{and}\quad
        \variance(X_1)=\sigma^2
    \]
    with \(0<\sigma^2<+\infty\). Then
    \[
        \frac{\sqrt{n}}{\sigma}\left(\overline{X}_n-\mu\right)
        \xrightarrow{d}
        Z,
    \]
    where \(Z\sim\operatorname{N}(0,1)\).
\end{snippettheorem}

\begin{snippetnote}{central-limit-theorem-normalization}{Normalization in the central limit theorem}
    Under the hypotheses of the \centrallimittheorem,
    \[
        S_n=\sum_{i=1}^{n}X_i
    \]
    satisfies
    \[
        \expectation[S_n]=n\mu
        \quad\text{and}\quad
        \variance(S_n)=n\sigma^2.
    \]
    Therefore
    \[
        \frac{\sqrt{n}}{\sigma}(\overline{X}_n-\mu)
        =
        \frac{S_n-n\mu}{\sigma\sqrt{n}}
    \]
    is the centered and normalized sum.
\end{snippetnote}

\begin{snippetexample}{gaussian-sample-mean-clt-example}{Gaussian sample means}
    Let \((X_n)_{n\geq1}\) be \iidrandomvariables such that
    \[
        X_n\sim\operatorname{N}(\mu,\sigma^2)
    \]
    with \(\sigma^2>0\). Then
    \[
        \overline{X}_n\sim\operatorname{N}\left(\mu,\frac{\sigma^2}{n}\right)
    \]
    and therefore
    \[
        \frac{\sqrt{n}}{\sigma}(\overline{X}_n-\mu)
        \sim
        \operatorname{N}(0,1)
    \]
    for every \(n\geq1\). In this case the normalized sample mean is
    already standard Gaussian for every \(n\), not only asymptotically.
\end{snippetexample}

\end{document}
